\begin{resumo}

A classificação de textos curtos em português, como as descrições de produtos de notas fiscais eletrônicas, enfrenta vocabulário abreviado e ruidoso e centenas de categorias desbalanceadas, e depende de grandes volumes de rótulos cuja obtenção é o principal custo desses sistemas. Esta tese propõe e avalia o FALCO (\textit{Framework} de Aprendizado Ativo com LLM para texto CurtO), que ataca esse custo em três frentes: composição informada do conjunto inicial de treinamento, seleção ativa de instâncias por incerteza e substituição parcial do anotador humano por oráculos baseados em modelos de linguagem de grande porte (LLMs), organizados em fases de custo crescente. A investigação usa um conjunto público auditado de 250.365 descrições de produtos em 621 categorias (250.221 na versão corrigida empregada nos experimentos) e segue protocolo pré-registrado com partições imutáveis, análise pareada (McNemar, Wilcoxon) e artefato rastreável para cada número reportado. Cinco resultados sustentam a tese. (i) A composição do conjunto inicial importa mais quanto menor o orçamento: só o sorteio produz amplitude de 6,4 pontos percentuais de acurácia com $|L_0|=100$, e a heurística determinística proposta, o DRI-SL, que combina densidade semântica e diversidade lexical sem usar nenhum rótulo, supera o melhor indivíduo de um algoritmo genético supervisionado em todos os tamanhos de 100 a 5.000; a comparação é conservadora, pois a correção de circularidade introduzida na tese reduziu o envelope evolutivo em até 6,3 pontos percentuais. (ii) Oráculos LLM \textit{zero-shot} atingem 77--83\% de acurácia no espaço fechado de 621 categorias, a custos de US\$~0,035 a US\$~0,92 por mil rótulos, com Macro F1 superior ao de um classificador supervisionado leve treinado com a base completa de 250 mil descrições rotuladas; quase metade dos seus erros concentra-se em convenções de catálogo (categorias-irmãs e classes guarda-chuva). (iii) O instrumento de medição é parte do resultado: sem restrição de saída ao espaço de classes, 6,8\% das respostas viram falsos erros; regras de fronteira no \textit{prompt} melhoram o modelo fraco na distribuição de produção ($+4{,}6$ p.p., $p=0{,}012$) ao custo de degradação severa nas classes raras ($p<0{,}001$); e o provedor de \textit{serving} também é instrumento, pois o mesmo modelo gratuito diverge significativamente de si mesmo entre dois serviços ($p<0{,}001$), embora, servido adequadamente, empate com oráculos pagos na amostra de produção. (iv) As estratégias de incerteza superam a seleção aleatória na significância máxima que 8 sementes permitem ($p=0{,}0078$), recuperando 78\% do Macro F1 do teto supervisionado com 15\% do \textit{pool} daquele experimento, e a vantagem sobrevive intacta a ruído uniforme de oráculo até $\varepsilon=0{,}4$, faixa que cobre o erro observado nos LLMs reais; como nenhum oráculo alcançou o critério pré-registrado de 85\% de acurácia, o gate metodológico definiu a configuração final (oráculo econômico na fase inicial, oráculo de maior acurácia na fase avançada) e tornou o aprendizado com rótulos ruidosos o cenário central. (v) O framework foi executado ponta a ponta com oráculo LLM gratuito, a custo monetário zero, e a validação com o classificador forte (BERTimbau, três sementes, avaliação na população reservada de 177.490 instâncias) respondeu à hipótese central, a de obter pelo menos $95\%$ da acurácia da supervisão completa do \textit{pool} de referência, a métrica do pré-registro, com no máximo 34.724 rótulos ($15\%$ da população deduplicada). O veredito: \textbf{o critério é atingível dentro do teto, cruzado com 20 mil rótulos ($8{,}6\%$ da população) nas três sementes, em varredura \textit{post hoc} com rótulos de gabarito, enquanto a configuração executada parou por estagnação aquém do piso, com 11.936 rótulos ($5{,}2\%$)}; a decomposição pareada inocenta o ruído do oráculo, que, por ser estruturado, até beneficia as classes raras, inocenta a seleção, e atribui a perda ao critério de parada herdado do classificador leve. Em Macro F1, análise de robustez introduzida nesta tese, o critério é atingido com 30 mil rótulos ($13{,}0\%$ da população), também nas três sementes e dentro do teto; e, na média das três sementes, 35 mil rótulos ativamente selecionados superam a supervisão completa do \textit{pool} nas duas métricas, com heterogeneidade declarada: duas sementes mantêm vantagem significativa, uma inverte em Macro F1 e empata em acurácia. Os experimentos em escala populacional mostram ainda que a autoavaliação em dados coletados ativamente é enviesada em direção que depende da métrica, pois a acurácia da autoavaliação subestima e o Macro F1 da autoavaliação superestima o desempenho real, o que fundamenta o conjunto reservado de avaliação e o critério de liberação que o framework adota. As limitações de primeira ordem são declaradas: os resultados provêm de uma única base, do próprio domínio de estudo; a superação da supervisão completa pelo braço de 35 mil rótulos vale na média das sementes, com uma semente divergente em Macro F1; e a análise de custo apoia-se em razões entre oráculos, como a diferença de 26 vezes sob empate estatístico de acurácia, e não em preços absolutos, que datam. A tese entrega ainda a métrica LCE (\textit{Learning Curve Efficiency}) para comparação de curvas de aprendizado, a biblioteca aberta \texttt{activelearning} e a ferramenta FlowBuilder, com ingestão saneada de dados, execução parametrizada e curvas de aprendizado, que tornam o protocolo reutilizável em outros domínios.

\end{resumo}
